\documentclass[10pt]{article}

\usepackage[margin=1in]{geometry}
\usepackage[T1]{fontenc}
\usepackage[utf8]{inputenc}
\usepackage{lmodern}
\usepackage{microtype}
\usepackage{amsmath,amssymb}
\usepackage{booktabs}
\usepackage{array}
\usepackage{longtable}
\usepackage{enumitem}
\usepackage{xcolor}
\usepackage{hyperref}
\usepackage{url}

\hypersetup{
  colorlinks=true,
  linkcolor=blue!55!black,
  citecolor=blue!55!black,
  urlcolor=blue!55!black,
  pdftitle={Libra: From Community Input to Audit Parameters},
  pdfauthor={Matthew Jackson}
}

\setlist{nosep,leftmargin=*}
\newcommand{\libra}{\textsc{Libra}}
\newcommand{\code}[1]{\texttt{#1}}

\title{\libra: From Community Input to Audit Parameters\\
\large A Provenance-Bearing Prototype for Participatory Fairness Auditing}
\author{Matthew Jackson\\Independent Researcher}
\date{September 30, 2026\\Version 1.0 preprint}

\begin{document}
\maketitle

\begin{abstract}
Fairness audits require choices about protected groups, reference populations,
metrics, and decision thresholds. Software libraries make these choices
configurable, but configuration authority usually remains with auditors,
developers, or regulators. This paper presents \libra, a research prototype
that represents an evaluation standard as a machine-readable configuration
containing priority groups, a reference group, a disparate-impact threshold,
evidence requirements, and provenance metadata. The prototype separates
researcher-default audits from audits backed by a documented community process
and implements a draft-to-finalized governance lifecycle for public-sector AI
evaluation artifacts.

We evaluate only the technical mechanism and threshold sensitivity; no
community participant data are used. On a deduplicated 2024 single-lender
Michigan HMDA extract ($n=3{,}827$), the Black-to-White observed origination
ratio is $0.8195$. On the ProPublica COMPAS dataset ($n=6{,}837$), the
African-American-to-Caucasian non-recidivism ratio is $0.8009$. Both point
estimates are above a $0.80$ threshold and below $0.85$, so their audit labels
change when the threshold changes. Bootstrap intervals cross these boundaries,
underscoring that threshold rules and statistical uncertainty must be reported
together. The result demonstrates parameter sensitivity, not that $0.85$ is
the correct threshold or that a community selected it. A publishable claim of
community validity requires prospective, ethically reviewed fieldwork.
\end{abstract}

\noindent\textbf{Keywords:} participatory AI; algorithmic auditing; disparate
impact; public-sector AI; provenance; community governance

\section{Introduction}

Algorithmic fairness is not reducible to selecting a single metric. Several
widely used fairness criteria are mutually incompatible except under restricted
conditions \cite{chouldechova2017,kleinberg2017}. Operational tools such as
AIF360, Aequitas, Fairlearn, and FairVis make measurement and mitigation more
accessible \cite{bellamy2018,saleiro2018,weerts2023,cabrera2019}, but a prior
governance question remains: who is authorized to choose the reference group,
the acceptable threshold, and the evidence sufficient to support a finding?

Participatory AI research has shown that affected people can contribute
substantive knowledge about system goals, harms, and trade-offs
\cite{lee2019,katell2020,birhane2022}. It has also warned that participation can
be extractive or symbolic when institutions retain control over scope and
implementation \cite{sloane2022}. Recent participatory-auditing studies further
show that non-specialists surface impacts absent from standard taxonomies, while
also needing scaffolding to translate lived concerns into measurable tests
\cite{zhao2026,dicampli2026,rezk2026,escher2026}.

This paper addresses a narrow systems problem within that larger literature:
how can the output of a documented community process be represented as a
versioned, provenance-bearing input to an audit of an existing public-sector AI
system? We present \libra, a prototype governance-to-audit interface. Its
contributions are:

\begin{enumerate}
  \item a formal representation of audit authority as a configuration
  $C=(G,r,\theta,P)$, where $G$ is a set of priority groups, $r$ is a reference
  group, $\theta$ is an audit threshold, and $P$ is provenance;
  \item an implemented prototype that records configuration source, supports a
  draft-to-finalized lifecycle, and carries the selected parameters into audit
  outputs; and
  \item a reproducible sensitivity analysis demonstrating that modest changes
  in $\theta$ alter which groups are flagged in two public benchmark datasets.
\end{enumerate}

The paper makes no claim that participatory auditing, provenance, registers, or
machine-readable AI governance are new categories. It does not report a
community-defined threshold, a completed civic pilot, government endorsement,
or improved social outcomes. The empirical contribution is a proof of
mechanism: parameter choice can change an audit classification, so authority
over that parameter is consequential.

\section{Related Work}

\subsection{Fairness metrics and audit toolkits}

Group-fairness research provides multiple, sometimes incompatible, ways to
compare outcomes and errors \cite{chouldechova2017,kleinberg2017}. AIF360,
Aequitas, Fairlearn, and FairVis provide metrics, visualization, and mitigation
methods \cite{bellamy2018,saleiro2018,weerts2023,cabrera2019}. These systems are
valuable technical foundations; \libra's focus is not metric invention. Its
focus is the authority and provenance attached to the parameters supplied to an
audit.

The four-fifths rule is an employment-selection heuristic in the U.S. Uniform
Guidelines, not a universal legal definition of discrimination
\cite{eeoc1978}. In this paper, $0.80$ is used as a familiar comparison scenario.
Applying that number to lending or criminal-justice data is illustrative and
does not establish legal compliance in either domain.

\subsection{Participation and community-informed auditing}

WeBuildAI demonstrated collective participation in designing an algorithmic
policy for food-assistance allocation \cite{lee2019}. The Algorithmic Equity
Toolkit supported community advocates auditing government technologies
\cite{katell2020}. Birhane et al. synthesize opportunities and unresolved power
questions in participatory AI, while Sloane et al. caution that participation
alone is not a design fix \cite{birhane2022,sloane2022}.

The closest recent work substantially narrows any broad novelty claim. Zhao et
al. co-designed community-centered auditing tools with 22 educators in Hawai'i
and argue that affected communities should shape who audits, the infrastructure
used, and the desired outcomes \cite{zhao2026}. Di Campli San Vito et al. show
that non-expert stakeholders identify overlooked harms but require assistance
operationalizing them \cite{dicampli2026}. Rezk et al. find that participatory
auditors' trust in apparently capable systems can reduce scrutiny
\cite{rezk2026}. Escher et al. conduct a community-informed audit grounded in
Kyrgyz caregivers' concerns about language and platform recommendations
\cite{escher2026}. These studies provide empirical participation evidence that
the present paper does not.

\subsection{Provenance and operational governance artifacts}

The Participation Ledger proposes machine-readable records linking community
contributions to consent, compensation, authority, system changes, and
replayable tests \cite{mushkani2026ledger}. Participatory provenance measures
how public submissions survive AI-mediated consultation summaries
\cite{mahajan2026}. Voices in the Loop offers a versioned and contestable atlas
of participatory-AI initiatives \cite{mushkani2026voices}.

Adjacent public-sector work also emphasizes reviewable artifacts. Vendor
factsheets can support dialogue but are insufficient as standalone risk
assessments \cite{kuehnert2026}. Li proposes a minimum reviewable trace for
AI-influenced public decisions \cite{li2026}. The Eticas taxonomy maps risk
concepts into tests, measurements, severity bands, and JSON-LD infrastructure
\cite{galdon2026}. A global comparison of 8,368 public-sector AI-register
records finds limited convergence and frequent omissions involving appeals,
risks, legal bases, and external evaluation \cite{das2026}. \libra is positioned
as a narrower prototype: a community-process output becomes an explicit audit
input whose origin remains visible in the result.

\section{Framework}

\subsection{Audit parameters}

Let $D=\{(x_i,a_i,y_i)\}_{i=1}^{n}$ be a dataset with sensitive-group label
$a_i\in A$ and binary favorable-outcome indicator $y_i\in\{0,1\}$. For group
$a$, the observed favorable-outcome rate is
\begin{equation}
  p_a = \Pr(Y=1\mid A=a)
      = \frac{\sum_i y_i\mathbb{1}[a_i=a]}
              {\sum_i \mathbb{1}[a_i=a]}.
\end{equation}
Given reference group $r$, the disparate-impact ratio is
\begin{equation}
  \operatorname{DI}(a,r)=\frac{p_a}{p_r}.
\end{equation}
A deterministic threshold rule flags group $a$ when
$\operatorname{DI}(a,r)<\theta$. Equality with the threshold does not trigger a
flag.

The configuration is
\begin{equation}
  C=(G,r,\theta,P,E),
\end{equation}
where $G\subseteq A$ contains priority groups, $P$ contains process provenance,
and $E$ contains evidence requirements such as minimum group size and acceptable
attestation sources. Table~\ref{tab:config} summarizes the prototype schema.

\begin{table}[ht]
\centering
\caption{Core configuration fields in the prototype.}
\label{tab:config}
\begin{tabular}{p{0.25\linewidth}p{0.67\linewidth}}
\toprule
Field & Purpose \\
\midrule
\code{priority\_groups} & Groups designated for elevated scrutiny. \\
\code{fairness\_target} & Reference group used to compute outcome ratios. \\
\code{fairness\_threshold} & Minimum accepted ratio under the configured rule. \\
\code{domain}, \code{jurisdiction} & Scope boundaries for reuse. \\
\code{provenance} & Process identifier, date, protocol, participant count,
facilitator, notes, and dissent or decision metadata where available. \\
\code{evidence\_requirements} & Minimum sample and evidence conditions applied
before substantive evaluation. \\
\code{attestor\_policy} & Permitted evidence signers or verifier classes. \\
\bottomrule
\end{tabular}
\end{table}

\subsection{Authority classification and lifecycle}

The prototype distinguishes a default-governed analysis from an analysis whose
configuration has passed a documented finalization workflow. This distinction
does not make a session representative merely because metadata fields are
present. ``Community-valid'' is an asserted governance classification whose
substantive validity depends on participant selection, informed consent,
facilitator independence, compensation, decision rules, dissent preservation,
scope, and authorization. Those conditions require field evidence and cannot
be proven by schema validation alone.

The API prototype supports draft session records, finalization, active-standard
resolution, provenance receipts, and audit outputs that identify whether a
default or finalized configuration was applied. The public registry is empty at
the time of this Version 1 analysis; therefore, the evaluation below produces
no community-valid audit.

\subsection{Implementation}

The implementation consists of a reusable Python package, a FastAPI service,
a React interface, a JSON-schema configuration format, and a controlled
registry. The audit engine computes group outcome rates, reference-group ratios,
threshold flags, and a max--min outcome-rate gap. The governance layer creates
and validates configuration objects and tracks lifecycle state. A separate
evidence path supports signed evidence and provenance receipts, but the present
experiment evaluates only the group-rate audit and configuration mechanism.

At the analyzed repository revision, 172 API and core tests passed. JSON
validation and both web builds also passed when run with the repository's Python
environment. These checks establish software behavior on the tested cases, not
social validity, legal compliance, production security, or deployment status.

\section{Evaluation}

\subsection{Research questions}

\begin{description}[leftmargin=3.6em,style=nextline]
  \item[RQ1] Does changing only the disparate-impact threshold alter which
  sufficiently represented groups are flagged?
  \item[RQ2] Can the reported result be regenerated from versioned code and
  local data artifacts without using participant data?
\end{description}

\subsection{Datasets}

\paragraph{HMDA Michigan.}
The source file contains 4,463 records from one 2024 Michigan lender filing.
Records with unavailable, free-form, joint, multi-minority, or ``Other'' race
labels were excluded. Four exact duplicate rows were removed, leaving 3,827
records. The favorable outcome is \code{action\_taken=1} (loan originated); all
other action codes are treated as non-origination. This is an observed
origination analysis, not an approval model, statewide estimate, or causal test
of discrimination. HMDA public data documentation is provided by the Consumer
Financial Protection Bureau \cite{cfpbHMDA}.

\paragraph{COMPAS.}
We use 7,214 records in ProPublica's Broward County COMPAS release
\cite{angwin2016}. Excluding the ``Other'' label leaves 6,837 records. The
favorable outcome is no recorded two-year recidivism. The dataset is a contested
historical benchmark and is used here only to reproduce a threshold-sensitivity
mechanism, not to endorse COMPAS or infer individual risk.

\subsection{Analysis}

We calculate group rates and ratios relative to White applicants in HMDA and
Caucasian defendants in COMPAS. We sweep
$\theta\in\{0.70,0.75,0.80,0.85,0.90,0.95\}$. Aggregate counts exclude groups
with fewer than 30 observations; excluded groups remain visible in the appendix.

To make sampling uncertainty explicit, we generate 10,000 stratified parametric
bootstrap replicates with fixed seed 20,260,930. Within each group, favorable
counts are sampled from a binomial distribution using that group's observed
rate. We report percentile intervals for the two focus ratios. These intervals
are descriptive and do not correct for dataset construction, measurement error,
or repeated analytic choices.

\section{Results}

\subsection{Threshold sensitivity}

Table~\ref{tab:sweep} reports threshold flags among groups with $n\geq30$.
Raising the scenario threshold from $0.80$ to $0.85$ adds the Black group in
HMDA and the African-American group in COMPAS. The Asian HMDA group is already
below $0.80$.

\begin{table}[ht]
\centering
\caption{Groups flagged at each threshold; groups with $n<30$ excluded from
counts.}
\label{tab:sweep}
\begin{tabular}{cccc}
\toprule
$\theta$ & HMDA & COMPAS & Combined \\
\midrule
0.70 & 0 & 0 & 0 \\
0.75 & 0 & 0 & 0 \\
0.80 & 1 & 0 & 1 \\
0.85 & 2 & 1 & 3 \\
0.90 & 2 & 1 & 3 \\
0.95 & 2 & 1 & 3 \\
\bottomrule
\end{tabular}
\end{table}

The two boundary cases are shown in Table~\ref{tab:focus}. Both point estimates
pass at $0.80$ and fail at $0.85$. Their intervals span one or both thresholds,
so a binary label should never be presented without the estimate, sample size,
and uncertainty.

\begin{table}[ht]
\centering
\caption{Focus-group results. Intervals are 95\% bootstrap percentile
intervals.}
\label{tab:focus}
\begin{tabular}{p{0.25\linewidth}rrrp{0.19\linewidth}}
\toprule
Dataset and group & $n_g$ & Rate & DI & 95\% interval \\
\midrule
HMDA: Black or African American & 630 & 0.3444 & 0.8195 & [0.7260, 0.9175] \\
COMPAS: African-American & 3,696 & 0.4857 & 0.8009 & [0.7652, 0.8391] \\
\bottomrule
\end{tabular}
\end{table}

These results answer RQ1 affirmatively as a deterministic property of the
configured rule: for a fixed point estimate, the threshold changes the label.
They do not show that $0.85$ is socially preferred, statistically optimal, or
legally required.

\subsection{Reproducibility}

The analysis script records dataset hashes, preprocessing counts, the repository
revision, all group rates, all threshold flags, the bootstrap seed, and the
number of replicates in a machine-readable result file. It can be regenerated
from the repository root with:

\begin{verbatim}
.venv/bin/python scripts/generate_preprint_v1_results.py
\end{verbatim}

At generation, the input SHA-256 hashes were
\code{17dad996...169d3029} for HMDA and
\code{c451db85...91d0cc7d} for COMPAS. A release archive and permanent artifact
identifier should be assigned before external submission.

\section{Discussion}

\subsection{What the experiment establishes}

The experiment establishes a narrow but consequential fact: an audit's
threshold is not a passive display setting. Holding the dataset, metric, and
reference group fixed, changing $\theta$ changes which groups receive a formal
flag. A governance process that supplies $\theta$ therefore influences the
audit's actionable output.

The result supports making parameter source explicit. A report should indicate
whether a threshold came from a regulation, an auditor default, a vendor, a
research scenario, or a documented community process. The same number has a
different authority claim under each source.

\subsection{What the experiment does not establish}

It does not establish that a community prefers $0.85$ or $0.90$. It does not
show that participants understand the consequences of the parameterization. It
does not validate the proposed facilitation protocol, demonstrate institutional
response, or show improved outcomes. It also does not turn a ratio threshold
into a legal conclusion. Those claims require different evidence.

\subsection{Design implications for prospective fieldwork}

Recent work suggests that a field study should go beyond asking participants to
move a numeric slider. Participants should help define harms, evidence sources,
perspectives, acceptable outcomes, and the authority to pause or reject a
system \cite{zhao2026,dicampli2026,mushkani2026ledger}. The process should record
disagreement instead of compressing it into false consensus, compensate
participants, report whose perspectives are absent, and specify who owns and
maintains the resulting artifact. ``Do not use this system'' must remain a valid
outcome. Otherwise, machine-readable participation can become participation
washing with better metadata.

\section{Limitations and Ethics}

\begin{itemize}
  \item \textbf{No participant study.} No residents, educators, public-sector
  staff, or affected decision subjects supplied data for this paper. No current
  configuration is described as community-defined.
  \item \textbf{Descriptive, non-causal analysis.} Group outcome ratios do not
  identify why disparities occur and cannot establish unlawful discrimination,
  model bias, or institutional intent.
  \item \textbf{Dataset scope.} The HMDA file covers one lender and treats every
  non-origination action as unfavorable. COMPAS is geographically and
  historically specific. Neither supports broad population generalization.
  \item \textbf{Categorization.} Administrative race labels flatten identity
  and omit intersectional structure. Some groups remain too small for stable
  interpretation; a minimum-$n$ rule reduces but does not solve this problem.
  \item \textbf{Uncertainty.} The bootstrap intervals cross decision
  thresholds. A mature implementation should incorporate interval-aware or
  evidentiary review instead of relying only on point-estimate flags.
  \item \textbf{Governance validity.} Provenance fields can document a process
  but cannot guarantee representativeness, independence, consent, or influence.
  Those properties require external review and institutional commitments.
  \item \textbf{Prototype maturity.} Passing tests does not establish
  production security or availability. The public deployment claim from an
  earlier draft was not verified for this paper and is intentionally omitted.
\end{itemize}

Any future study that collects participant responses for publication should
receive an institutional ethics or IRB determination before recruitment. It
should use informed consent, privacy minimization, withdrawal procedures,
participant compensation, and explicit authorization for quotations and
artifact publication. District 3 and the City of Detroit are prospective
contexts only; neither is represented as an endorser, adopter, or completed
pilot in this paper.

\section{Conclusion}

\libra demonstrates a governance-to-audit mechanism in which the source and
content of audit parameters remain visible in machine-readable outputs. Across
two public benchmark datasets, moving a disparate-impact threshold from $0.80$
to $0.85$ changes the status of substantively large groups. This shows why the
choice and authority of an audit standard should be documented. Version 1 is a
technical and methodological proof of mechanism. The next evidentiary step is a
prospective, ethically reviewed co-design study that tests whether the protocol
produces an authorized and useful standard and whether institutions act on it.

\appendix
\section{Group-Level Descriptive Results}

\begin{longtable}{p{0.39\linewidth}rrrcl}
\caption{All analyzed groups. A dagger marks groups excluded from aggregate
flag counts because $n<30$.}\label{tab:allgroups}\\
\toprule
Dataset and group & $n$ & Rate & DI & $<0.80$ & $<0.85$ \\
\midrule
\endfirsthead
\toprule
Dataset and group & $n$ & Rate & DI & $<0.80$ & $<0.85$ \\
\midrule
\endhead
HMDA: American Indian or Alaska Native$^{\dagger}$ & 29 & 0.1379 & 0.3281 & yes & yes \\
HMDA: Asian & 232 & 0.3233 & 0.7691 & yes & yes \\
HMDA: Black or African American & 630 & 0.3444 & 0.8195 & no & yes \\
HMDA: Native Hawaiian or Other Pacific Islander$^{\dagger}$ & 5 & 0.2000 & 0.4758 & yes & yes \\
HMDA: White (reference) & 2,931 & 0.4203 & 1.0000 & no & no \\
\midrule
COMPAS: African-American & 3,696 & 0.4857 & 0.8009 & no & yes \\
COMPAS: Asian & 32 & 0.7188 & 1.1854 & no & no \\
COMPAS: Caucasian (reference) & 2,454 & 0.6064 & 1.0000 & no & no \\
COMPAS: Hispanic & 637 & 0.6358 & 1.0485 & no & no \\
COMPAS: Native American$^{\dagger}$ & 18 & 0.4444 & 0.7330 & yes & yes \\
\bottomrule
\end{longtable}

\section*{Data and Code Availability}

The Version 1 analysis code, machine-readable results, and manuscript source
are available in the versioned public release at
\url{https://github.com/MattJaxson/libra-preprint-v1/releases/tag/v1.0.0}.
The manuscript is licensed under Creative Commons Attribution 4.0
International (CC BY 4.0), and original analysis code is licensed under the
Apache License 2.0. The raw HMDA and COMPAS files derive from third-party public
sources, are not relicensed by this paper, and should be obtained and used in
accordance with their source terms. The manuscript contains no confidential
District 3 materials and no participant-level research data.

\begin{thebibliography}{99}

\bibitem{angwin2016}
J. Angwin, J. Larson, S. Mattu, and L. Kirchner.
Machine bias.
\emph{ProPublica}, May 23, 2016.
\url{https://www.propublica.org/article/machine-bias-risk-assessments-in-criminal-sentencing}

\bibitem{bellamy2018}
R. K. E. Bellamy et al.
AI Fairness 360: An extensible toolkit for detecting, understanding, and
mitigating unwanted algorithmic bias.
\emph{arXiv:1810.01943}, 2018.

\bibitem{birhane2022}
A. Birhane, W. Isaac, V. Prabhakaran, M. D\'iaz, M. C. Elish, I. Gabriel,
and S. Mohamed.
Power to the people? Opportunities and challenges for participatory AI.
\emph{EAAMO}, 2022. \url{https://arxiv.org/abs/2209.07572}

\bibitem{cabrera2019}
\'A. A. Cabrera, W. Epperson, F. Hohman, M. Kahng, J. Morgenstern, and
D. H. Chau.
FairVis: Visual analytics for discovering intersectional bias in machine
learning.
\emph{arXiv:1904.05419}, 2019.

\bibitem{cfpbHMDA}
Consumer Financial Protection Bureau.
Home Mortgage Disclosure Act data.
\url{https://ffiec.cfpb.gov/data-publication/}

\bibitem{chouldechova2017}
A. Chouldechova.
Fair prediction with disparate impact: A study of bias in recidivism prediction
instruments.
\emph{Big Data}, 5(2):153--163, 2017.
doi:10.1089/big.2016.0047.

\bibitem{das2026}
D. Das and S. Guha.
A global comparison of schemas, transparency, and interoperability in
public-sector AI registers and inventories.
\emph{arXiv:2609.24883}, 2026.

\bibitem{dicampli2026}
P. Di Campli San Vito et al.
Empowering stakeholders with participatory auditing of predictive AI:
Perspectives from end-users and decision subjects without AI expertise.
In \emph{Proceedings of CHI 2026}, pages 1--35, 2026.
doi:10.1145/3772318.3791757.

\bibitem{eeoc1978}
U.S. Equal Employment Opportunity Commission.
Uniform Guidelines on Employee Selection Procedures, 29 C.F.R.
\S~1607.4(D), 1978.

\bibitem{escher2026}
N. Escher, B. Yrysov, A. McDermott, D. Chechelnitsky, H. B. Benyam, and
N. Banovic.
Empire Amplifier: Uncovering and contesting the prioritization of colonial
content on platforms through community-informed algorithmic auditing.
\emph{arXiv:2604.27498}, 2026.

\bibitem{galdon2026}
G. Galdon Clavell, P. Accuosto, and U. Gohar.
The Eticas AI Risk Taxonomy: Open infrastructure for operationalizing AI
audits.
\emph{arXiv:2607.02201}, 2026.

\bibitem{katell2020}
M. Katell, M. Young, B. Herman, D. Dailey, A. Tam, V. Guetler, C. Binz,
D. Raz, and P. M. Krafft.
An algorithmic equity toolkit for technology audits by community advocates and
activists.
In \emph{Proceedings of FAccT 2020}, 2020.
\url{https://arxiv.org/abs/1912.02943}

\bibitem{kleinberg2017}
J. Kleinberg, S. Mullainathan, and M. Raghavan.
Inherent trade-offs in the fair determination of risk scores.
In \emph{Proceedings of ITCS 2017}, 2017.
\url{https://arxiv.org/abs/1609.05807}

\bibitem{kuehnert2026}
B. Kuehnert, N. Johnson, R. Dotan, and H. Heidari.
Disclosure or marketing? Analyzing the efficacy of vendor self-reports for
vetting public-sector AI.
\emph{arXiv:2604.01332}, 2026.

\bibitem{lee2019}
M. K. Lee, D. Kusbit, A. Kahng, J. T. Kim, X. Yuan, A. Chan, D. See,
R. Noothigattu, S. Lee, A. Psomas, and A. D. Procaccia.
WeBuildAI: Participatory framework for algorithmic governance.
\emph{Proceedings of the ACM on Human-Computer Interaction}, 3(CSCW), 2019.
doi:10.1145/3359283.

\bibitem{li2026}
G. Li.
Auditable accountability without an AI Act: Australia's public-sector AI
assurance stack and the minimum reviewable trace.
\emph{Law, Ethics \& Technology}, 3(3):1--13, 2026.
doi:10.55092/let20260010.

\bibitem{mahajan2026}
S. Mahajan.
Participatory provenance as representational auditing for AI-mediated public
consultation.
\emph{arXiv:2604.20711}, 2026.

\bibitem{mushkani2026ledger}
R. Mushkani.
Traceable, enforceable, and compensable participation: A Participation Ledger
for people-centered AI governance.
\emph{arXiv:2602.10916}, 2026.

\bibitem{mushkani2026voices}
R. Mushkani.
Voices in the Loop: Mapping participatory AI.
\emph{Proceedings of FAccT 2026}, 2026.
\url{https://arxiv.org/abs/2605.16827}

\bibitem{rezk2026}
A. M. Rezk, P. Di Campli San Vito, A. Soufan, G. McDonald, C. Macdonald,
and I. Ounis.
All Eyes on the Ranker: Participatory auditing to surface blind spots in ranked
search results.
\emph{arXiv:2604.09946}, 2026.

\bibitem{saleiro2018}
P. Saleiro, B. Kuester, L. Hinkson, J. London, A. Stevens, A. Anisfeld,
K. T. Rodolfa, and R. Ghani.
Aequitas: A bias and fairness audit toolkit.
\emph{arXiv:1811.05577}, 2018.

\bibitem{sloane2022}
M. Sloane, E. Moss, O. Awomolo, and L. Forlano.
Participation is not a design fix for machine learning.
In \emph{Proceedings of EAAMO}, 2022.
\url{https://arxiv.org/abs/2007.02423}

\bibitem{weerts2023}
H. Weerts, M. Dud\'ik, R. Edgar, A. Jalali, R. Lutz, and M. Madaio.
Fairlearn: Assessing and improving fairness of AI systems.
\emph{Journal of Machine Learning Research}, 24(257):1--8, 2023.

\bibitem{zhao2026}
D. Zhao, H. Cha, M. J. Ryan, A. Wang, R. Baker-Ramos,
E.-B. Helekahi-Kaiwi, R. Diego, J. Hester, and D. Yang.
Whose Knowledge Counts? Co-designing community-centered AI auditing tools with
educators in Hawai'i.
In \emph{Proceedings of CHI 2026}, pages 1--24, 2026.
doi:10.1145/3772318.3790958.

\end{thebibliography}

\end{document}
